\section{Charging the solver at the true size of the middle part}
\label{sec:middle}

In the reduction of Section~\ref{sec:host} the middle part of an instance is
\(C_k\times\Z_p\) with \(|C_k|\le\lceil s/g'\rceil\) and \(p\le s\),
\(s=\lfloor\sqrt{D'}\rfloor\).  So it has at most
\[
  \operatorname{mid}(D',g'):=\Bigl\lceil\frac{s}{g'}\Bigr\rceil s\approx\frac{D'}{g'}
\]
vertices.  The paper hands the instance to the solver as an instance of
\(\LopTri(n,D')\) and charges it accordingly.  Handing it over as an instance
of \(\LopTri(n,\operatorname{mid}(D',g'))\) changes one assignment of the
program and separates two roles of \(D'\): the number of query pairs is
governed by \(\sqrt{D'}\), the cost of a call by \(D'/g'\).

\subsection{The reduction}

\begin{theorem}[Theorem~17 at the true middle size]\label{thm:mid}
There is a constant \(C\) with the following property.  Let \(T(n,D,w)\) bound
the time of a solver of \(\LopTri(n,D)\) with \(w\) query pairs.  Then Exact
Triangle on \(n\) vertices per part with weights of absolute value at most
\(n^{\nu}\), \(\nu\ge1\), is solved in time at most
\[
  4ng'\Bigl(T\bigl(n,\operatorname{mid}(D',g'),\lfloor n^{2}/\sqrt{D'}\rfloor\bigr)
     +C\frac{n^{2}}{\sqrt{D'}}\Bigr)
  +C\Bigl(\frac{\nu n^{3}\log n}{g'}+n^{\log_2 7}D'^{3/2}+n^{2}D'g'\Bigr)
\]
whenever \(16\le D'\le n\) and \(1\le g'\le\sqrt{D'}\).
\end{theorem}

This is a statement about programs: the solver is a procedure of a program of
the machine, and the theorem gives a program for Exact Triangle.  The exponent
\(\log_2 7\) is that of Strassen's algorithm \cite{Str69}, which the upstream
program uses for the choice of the prime; see Section~\ref{sec:limits}.

In Lean the statement is \lean{Claim17Mid}, the upstream statement
\lean{Claim.Theorem_17} with \(\operatorname{mid}(D',g')\) (\lean{midSize}) in
place of \(D'\) in the first argument of \(T\).  The host is upstream's with one
added assignment (\lean{et17Sizes'}): after the prime and the size of a piece
are computed, the parameter that is passed to the solver is overwritten with
\(\lceil s/g'\rceil s\).  Its correctness and its time follow from upstream's by
rewriting the arguments of the solver (\lean{hostTime'_eq},
\lean{obeysBound17Mid_hostTime}); the claim for programs is
\lean{claim17Mid_of_paramProcs}, and for \(D'=\lfloor n^{a/b}\rfloor\),
\(g'=\lceil D'^{c/d}\rceil\) with rational exponents it is
\lean{hostBound_mid_rat}.

\subsection{The total time}

\begin{theorem}[Total time]\label{thm:total}
Let \(\gamma,q\ge0\) and \(\varepsilon>0\).  Suppose that \(\LopTri(n,D)\) with
\(w\) query pairs is solved, for \(D\le n^{\varepsilon}\), in time
\(O\bigl(n^{2}(\log D+1)^{2}/D^{\gamma}+w\,D^{q}(\log D+1)\bigr)\).  Let
\(0<d\) and \(0<\eta<1/2\) be rational with \(d(1-\eta)<\varepsilon\), and let
\(D'=\lfloor n^{d}\rfloor\), \(g'=\lceil D'^{\eta}\rceil\).  Then Exact Triangle
with weights up to \(n^{\nu}\) is solved in time \(O(\nu\,n^{3-\delta}\log n)\)
for every \(\delta\ge0\) with
\begin{equation}\label{eq:total}
  \delta<d\,\min\Bigl(\eta,\ \min\bigl(\gamma(1-\eta),\ \tfrac12-q(1-\eta)\bigr)-\eta\Bigr)
  \qquad\text{and}\qquad \tfrac32d+\delta\le0.19 .
\end{equation}
\end{theorem}

\begin{proof}
Apply Theorem~\ref{thm:mid}.  The dimension of a call is
\(D=\operatorname{mid}(D',g')=\Theta(D'^{1-\eta})=\Theta(n^{d(1-\eta)})\), which
is at most \(n^{\varepsilon}\) for large \(n\).  Up to constants and powers of
\(\log n\), the terms are
\begin{align*}
  &\text{calls, preprocessing:} & n\,g'\,n^{2}D^{-\gamma}&=n^{3+d\eta-d\gamma(1-\eta)},\\
  &\text{calls, queries:} & n\,g'\,\frac{n^{2}}{\sqrt{D'}}D^{q}&=n^{3+d\eta+dq(1-\eta)-d/2},\\
  &\text{scans:} & \nu n^{3}/g'&=\nu n^{3-d\eta},\\
  &\text{prime:} & n^{\log_2 7}D'^{3/2}&\le n^{2.81+3d/2},\\
  &\text{construction:} & n^{2}D'g'&=n^{2+d+d\eta}.
\end{align*}
The first three exponents are at most \(3-\delta\) by the first condition of
\eqref{eq:total}, with room for the logarithms, and so is that of the term
\(ng'n^{2}/\sqrt{D'}\).  The last two are at most \(3-\delta\) by the second
condition.
\end{proof}

In Lean: \lean{explicitFrom_mid}, from \lean{hostBound_mid_rat} and a solver
of \(\LopTri\) (\lean{LopClaim}, obtained from a thin-product claim by
\lean{lopClaim_of_thinClaim}).  The same calculation with every call charged at
\(D'\), as in the paper, is \lean{explicitFrom_original}; its condition is
\(\delta<d\min(\eta,\min(\gamma,\frac12-q)-\eta)\).

\subsection{The saving of a tuple}

It is convenient to take the dimension \(D\) of the calls as the basic
parameter.  Write \(D\approx n^{\varepsilon}\) and
\begin{equation}\label{eq:sigma}
  \sigma=\frac{1}{2(1-\eta)}\in\bigl[\tfrac12,1\bigr),\qquad\text{so that}\qquad
  D'=D^{2\sigma},\quad g'=D^{2\sigma-1},\quad\sqrt{D'}=D^{\sigma}.
\end{equation}
An instance then has inner dimension \(D\) and at most \(n^{2}/D^{\sigma}\)
query pairs.  With the thin product of Section~\ref{sec:exact} at the
parameters \((c,\theta)\) the exponent of a call is
\(\min(\gX(c,\theta),\sigma-q(\theta))\), by \eqref{eq:cor32} with
\(\kappa=\sigma\), and the first condition of \eqref{eq:total} reads
\(\delta<\sav(c,\theta,\sigma,\varepsilon)\) with \(d(1-\eta)\) in the place of
\(\varepsilon\), where:

\begin{definition}\label{def:saving}
The \emph{saving} of the tuple \((c,\theta,\sigma,\varepsilon)\) is
\[
  \sav(c,\theta,\sigma,\varepsilon)
  =\varepsilon\,\min\Bigl(2\sigma-1,\ \min\bigl(\gX(c,\theta),\sigma-q(\theta)\bigr)-(2\sigma-1)\Bigr).
\]
\end{definition}

In Lean: \lean{tupleSaving}, with \lean{effGamma} for the inner minimum.  The
first argument of the outer minimum is the saving of the scans, the second
that of the calls.  In the paper every call is charged at the dimension \(D'\); relative to
\(D'\) the density of the query pairs is \(\kappa=1/2\).  Relative to the true
dimension \(D\) it is \(\sigma>1/2\) as soon as \(g'>1\), and the thin product
profits from both: a smaller inner dimension and, relative to it, fewer wanted
entries.

\begin{remark}[A larger hash prime]\label{rem:prime}
The same algorithm can be described without \(g'\).  Hash modulo a prime
\(p\approx D^{\sigma}\) with \(\sigma\ge1/2\), and cut each part into pieces of
\(D^{1-\sigma}\) vertices, so that a middle part \(C_k\times\Z_p\) has \(D\)
vertices.  There are \(O(nD^{2\sigma-1})\) instances with at most
\(n^{2}/D^{\sigma}\) query pairs each, and the scans cost
\(\nu n^{3}\log n/D^{2\sigma-1}\).  This is the reduction above with
\(D'=D^{2\sigma}\) and \(g'=D^{2\sigma-1}\): a larger prime and no further
splitting.  The balance of Remark~20 of the paper, \(\eta=\gamma/2\), becomes
the fixed point \(\sigma=1/2+\gamma/4\) of Section~\ref{sec:optimum}.
\end{remark}
